% --- Archon physics formalization source begin ---
% archon:physics
% archon:covers IPhO2026Problems/problem_IPhO_2026_2_C_3.lean
% archon:source-report reports/ipho_2026_k3/problem_IPhO_2026_2_C_3.source.json
% archon:problem-id IPhO_2026_2
% archon:part-id C.3
% archon:previous-part IPhO_2026_2_C_1
% archon:previous-part-policy natural_language_prerequisite_only
% archon:previous-part IPhO_2026_2_C_2
% archon:previous-part-policy natural_language_prerequisite_only

\chapter{Physics problem IPhO\_2026\_2\_C\_3}
\label{ch:IPhO2026Problems_problem_IPhO_2026_2_C_3}

\paragraph{Problem source.}
For the half-cylindrical mirror of radius R, ray A is incident at angle theta
and its reflected line is y = m\_A*x + b\_A.  A neighboring parallel ray B is
incident at theta + Delta theta, with Delta theta much smaller than theta, and
its reflected line is y = m\_B*x + b\_B.  The envelope/intersection of neighboring
rays forms the caustic.  Use Figure 2g and its coordinate convention.

Current subquestion:
Find the limiting intersection coordinates (X\_c,Y\_c) of the neighboring reflected rays.

\paragraph{Current subquestion.}
Find the limiting intersection coordinates (X\_c,Y\_c) of the neighboring reflected rays.

\paragraph{Recorded answer/context.}
X\_c = R*sin(theta)\textasciicircum{}3; Y\_c = (R/2)*cos(theta)*(2 - cos(2*theta)).

\paragraph{Figure/image path.}
/root/proposal\_for\_physic/science-mango/ipho\_2026\_source/image/T2\_page-4.png

\paragraph{Reusable previous-part conclusions.}
\begin{itemize}
\item Source C.1. Question: Write the slope m\_A and intercept b\_A of reflected ray A in terms of theta and R. Reusable conclusions: m\_A = cot(2*theta), and b\_A = R/(2*cos(theta)). Policy: natural\_language\_prerequisite\_only; do\_not\_import\_Lean\_output
\item Source C.2. Question: Expand m\_B and b\_B to first order in Delta theta. Reusable conclusions: m\_B = cot(2*theta) - 2*csc(2*theta)\textasciicircum{}2*Delta theta; b\_B = [R/(2*cos(theta))]*(1 + tan(theta)*Delta theta), up to O(Delta theta\textasciicircum{}2). Policy: natural\_language\_prerequisite\_only; do\_not\_import\_Lean\_output
\end{itemize}

\paragraph{Formalization target.}
create a compiling Lean file with sorry bodies at `IPhO2026Problems/problem\_IPhO\_2026\_2\_C\_3.lean`.
The Lean declarations must preserve the physical quantities, dimensions or dimensional roles, figure labels, governing-law hypotheses, and final relation expressed by this problem.
Use Mathlib/Physlib names found through LeanExplore where available. If a domain API is missing, introduce faithful local abstractions rather than scalar placeholder aliases.

\begin{theorem}[Physics formalization target]
\label{thm:physics:IPhO_2026_2_C_3:target}
The assigned autoformalize agent should translate this physics problem into Lean declarations in the covered file, with theorem and lemma proof bodies written as `by sorry`.
\end{theorem}
\begin{proof}
This is an autoformalization task, not a proof task. Produce faithful statements that can later be proved without weakening the source contract.
\end{proof}
% --- Archon physics formalization source end ---

% --- Archon autoformalization link (added by autoformalize agent) ---
\begin{theorem}[Limiting intersection of neighboring reflected rays]
\label{thm:IPhO2026Problems_problem_IPhO_2026_2_C_3:limitingIntersectionCoordinates}
\lean{IPhO2026Problems.IPhO2026_2_C_3.limitingIntersectionCoordinates}
\leanok
\uses{thm:physics:IPhO_2026_2_C_3:target}
For the half-cylindrical mirror of radius $R$ in the coordinate system of
Figure~2g, the intersection of the reflected ray A (incidence angle $\theta$)
with the reflected neighboring ray B (incidence angle $\theta+\Delta\theta$)
tends, as $\Delta\theta\to 0$, to the caustic point
\[
X_c = R\sin^3\theta, \qquad Y_c = \frac{R}{2}\cos\theta\,\bigl(2-\cos 2\theta\bigr).
\]
\end{theorem}
\begin{proof}
Formalized in Lean with a `sorry` body; the proof solves the two line
equations for the intersection, divides by $\Delta\theta$, and uses the
first-order expansions of part C.2 together with the part C.1 values of
$m_A$ and $b_A$.
\end{proof}
